\chapter{Аффинное дефектное дополнение: no-go размерностного совпадения}
\label{ch:v7-affine-defect-bicomplex-completion-gate}

\section{Проверяемая зацепка}

Предыдущий гейт обнаружил равенство
\begin{equation}
 54-42=12=\dim_{\mathbb C}E_{\rm aff}.
 \label{eq:v7-affine-defect-apparent-match}
\end{equation}
Возникла возможность дополнить связывающий носитель аффинным модулем и
получить две изометричные $54$-мерные стороны. Но такое заключение допустимо
только в том случае, если числа $54$ и $42$ являются инвариантными
размерностями минимальных кривизностных опор, а предполагаемый дефект несёт
то же действие алгебры, что и $E_{\rm aff}$.

\section{Тождественно нулевая средняя ступень}

Связывающая цепь имеет вид
\begin{equation}
 H_0\longrightarrow H_1\longrightarrow H_2,
 \qquad
 (\dim H_0,\dim H_1,\dim H_2)=(11,21,10).
\end{equation}
Её полная размерность равна $42$. Однако относительная кривизна
\begin{equation}
 \mathbb M_L=i\,\delta_N(Q_A^2-Q_{A_0}^2)
\end{equation}
содержит только крайние блоки
\begin{equation}
 \mathcal R_U:H_0\to H_2,
 \qquad
 \mathcal R_U^*:H_2\to H_0.
\end{equation}
На всей средней ступени $H_1$ она тождественно равна нулю:
\begin{equation}
 P_{H_1}\mathbb M_L=0,
 \qquad
 \mathbb M_LP_{H_1}=0.
 \label{eq:v7-affine-defect-middle-zero}
\end{equation}
Сто случайных физических конфигураций дали нулевой машинный остаток.

Если $J_{02}:H_0\oplus H_2\hookrightarrow H_0\oplus H_1\oplus H_2$ есть
каноническое включение концов, то
\begin{equation}
 \widetilde M_L=J_{02}^*\mathbb M_LJ_{02}in M_{21}(\mathbb C)
\end{equation}
удовлетворяет точному тождеству
\begin{equation}
 \frac12\Tr_{42}\mathbb M_L^2
 =\frac12\Tr_{21}\widetilde M_L^2.
 \label{eq:v7-affine-defect-support-compression}
\end{equation}
Максимальная невязка меньше $2.3\times10^{-13}$. На текущем физическом
срезе максимальный ранг обеих матриц равен $14$.

\section{Нулевая добавка меняет предполагаемый дефект}

Для любого $k\ge0$ оператор
\begin{equation}
 \mathbb M_L^{(k)}=\mathbb M_L\oplus0_k
\end{equation}
даёт то же действие:
\begin{equation}
 \Tr(\mathbb M_L^{(k)})^2=\Tr\mathbb M_L^2.
 \label{eq:v7-affine-defect-zero-padding}
\end{equation}
Проверены размерности $42$, $54$ и $79$ с одинаковым следом. Поэтому
разность $54-42=12$ зависит от выбора нулевого дополнения. После сжатия к
канонической endpoint-опоре она превращается в
\begin{equation}
 54-21=33.
 \label{eq:v7-affine-defect-support-difference}
\end{equation}
Следовательно, число $12$ не является индексом, рангом кривизности или
размерностью канонического quotient.

\section{Бимодульная несовместимость}

Аффинный модуль имеет конкретное действие $S_4$:
\begin{equation}
 E_{\rm aff}=\operatorname{Hom}(\mathbb C^4,\operatorname{im}P_3),
 \qquad X\longmapsto r(g)X\pi(g)^*.
\end{equation}
Его точное разложение на неприводимые представления равно
\begin{equation}
 E_{\rm aff}\simeq
 \mathbf1\oplus\mathbf2\oplus2\,\mathbf3\oplus\mathbf3'.
 \label{eq:v7-affine-defect-s4-decomposition}
\end{equation}
В частности, инвариантное подпространство одномерно.

У предполагаемого дефекта внутри искусственного $54$-мерного контейнера
действие $S_4$ вообще не задано. Если использовать единственный
factor-blind выбор, то есть тривиальное действие, любой эквивариантный
оператор из $E_{\rm aff}$ имеет ранг не выше единицы. Машинное усреднение
пятидесяти общих матриц подтвердило этот предел с остатком меньше
$2.1\times10^{-15}$. Полный изоморфизм можно получить только если отдельно
перенести на дефект всё представление \eqref{eq:v7-affine-defect-s4-decomposition};
это и было бы новым недоказанным входом.

Такое различение согласуется с классификацией конечных спектральных троек:
структура конечного бимодуля определяется матрицами кратностей и
дополнительными Real/градуированными данными, а не одной полной
размерностью \cite{v7-affine-defect-krajewski,v7-affine-defect-cacic}.

Наконец, даже выбор вложения $\mathbb C^{42}\subset\mathbb C^{54}$ не
каноничен. Пространство таких $42$-мерных подпространств есть
\begin{equation}
 \operatorname{Gr}_{\mathbb C}(42,54),
 \qquad
 \dim_{\mathbb R}=2\cdot42\cdot12=1008.
 \label{eq:v7-affine-defect-grassmannian}
\end{equation}
Без дополнительной симметрии предполагаемый дефект является выбором среди
1008 вещественных параметров.

\begin{theorem}[No-go аффинного дефектного дополнения]
Равенство $54-42=12=\dim_{\mathbb C}E_{\rm aff}$ не является инвариантным
совпадением физических носителей. Размерность $42$ содержит тождественно
нулевую среднюю ступень и может произвольно меняться нулевым дополнением;
после канонического сжатия активный endpoint-носитель имеет размерность
$21$. Кроме того, никакое действие алгебры на предполагаемом дефекте не
выведено. Поэтому каноническое отображение из $E_{\rm aff}$ отсутствует.
\end{theorem}

\begin{remark}
Этот запрет не затрагивает общий след и качественный вакуум предыдущих
гейтов. Он удаляет только ложную надежду вывести количественную метрику из
разности размеров двух выбранных матричных контейнеров.
\end{remark}

\begin{nogocriterion}[Запрет дополнения по одному числу]
Нельзя добавлять $E_{\rm aff}$ к связывающему носителю только потому, что
$54-42=12$. Дополнение допустимо лишь после предъявления общего действия
алгебры, эквивариантного отображения полного ранга и независимого
физического смысла нулевой средней ступени.
\end{nogocriterion}

\begin{protocol}
\begin{enumerate}
 \item полная размерность цепи: \textbf{$42=11+21+10$};
 \item средняя ступень в relative curvature: \textbf{тождественный нуль};
 \item минимальный endpoint-носитель: \textbf{$21=11+10$};
 \item сохранение следа после сжатия: \textbf{получено};
 \item максимальный физический ранг кривизны: \textbf{$14$};
 \item инвариантность при нулевом дополнении: \textbf{получена};
 \item предполагаемый дефект $54-42$: \textbf{$12$, неинвариантен};
 \item разность с минимальной опорой: \textbf{$33$};
 \item $S_4$-инвариантных линий в $E_{\rm aff}$: \textbf{$1$};
 \item эквивариантный изоморфизм дефекта: \textbf{не получен};
 \item свобода выбора $42$-плоскости: \textbf{$1008$ вещественных параметров};
 \item качественный вакуум: \textbf{не изменён};
 \item массовые отношения: \textbf{не выведены};
 \item следующий гейт: \textbf{минимальная опора кривизны и её след}.
\end{enumerate}
\end{protocol}

Машинный сертификат сохранён в
\path{s2t_v7_affine_defect_bicomplex_completion_gate_results.json}.

\begin{thebibliography}{9}
\bibitem{v7-affine-defect-krajewski}
T.~Krajewski,
\emph{Classification of Finite Spectral Triples},
J. Geom. Phys. 28 (1998), 1--30; arXiv:hep-th/9701081.
\bibitem{v7-affine-defect-cacic}
B.~\v{C}a\'ci\'c,
\emph{Moduli Spaces of Dirac Operators for Finite Spectral Triples},
J. Noncommut. Geom. 3 (2009), 453--495; arXiv:0902.2068.
\end{thebibliography}